% --- Archon physics formalization source begin ---
% archon:physics
% archon:covers PhyXMiniProblems/problem_phyx_mini_0857.lean
% archon:source-report reports/phyx_mini/problem_phyx_mini_0857.source.json

\chapter{Physics problem phyx\_mini\_0857}
\label{ch:PhyXMiniProblems_problem_phyx_mini_0857}

\paragraph{Problem source.}
\#\# Question

What is the electric potential at point C?

\#\# Answer choices

- A: A: 1100V

- B: B: 1600V

- C: C: 800V

- D: D: 900V

\#\# Figure caption (auxiliary; use the image as primary evidence)

The image depicts a diagram of electric field lines surrounding a positively charged particle located at the center. The positive charge is marked as "2.0 nC."

Three points labeled A, B, and C are shown at different locations around the charge. Dotted circles denote equidistant lines from the charge. 

- Point B is located at a distance of 1.0 cm from the charge.
- Point A is positioned outside the 1.0 cm circle.
- Point C lies outside the 2.0 cm circle.

An arrow indicates that point B is 1.0 cm from the charge, and another arrow shows point C at 2.0 cm. The setup illustrates the spatial relationship between the charge and the points A, B, and C.

\#\# Dataset metadata

Electrostatics; Physical Model Grounding Reasoning, Spatial Relation Reasoning

\paragraph{Recorded answer/context.}
D: D: 900V

\paragraph{Figure/image path.}
/root/proposal\_for\_physic/science-mango/phyx\_mini\_run/phyx\_data/test\_image/857.png

\paragraph{Formalization target.}
create a compiling Lean file with sorry bodies at `PhyXMiniProblems/problem\_phyx\_mini\_0857.lean`. The Lean declarations must preserve the physical quantities, dimensions or dimensional roles, figure labels, governing-law hypotheses, and final relation expressed by this problem.
Use Mathlib/Physlib names found through LeanExplore where available. If a physics API is missing, introduce faithful local abstractions rather than scalar placeholder aliases.

\begin{theorem}[Physics formalization target]
\label{thm:physics:phyx_mini_0857:target}
\lean{PhyXMiniProblems.ProblemPhyXMini0857.problem_phyx_mini_0857}
\uses{def:physics:phyx-mini-0857:phyxminiproblems-problemphyxmini0857-electricpotentialinvolts, def:physics:phyx-mini-0857:phyxminiproblems-problemphyxmini0857-pointchargepotentialsetup, def:physics:phyx-mini-0857:phyxminiproblems-problemphyxmini0857-matchessuppliedpointchargefigure, def:physics:phyx-mini-0857:phyxminiproblems-problemphyxmini0857-usesschoolcoulombconstant, def:physics:phyx-mini-0857:phyxminiproblems-problemphyxmini0857-hasphysicalpointchargeparameters, def:physics:phyx-mini-0857:phyxminiproblems-problemphyxmini0857-satisfiespointchargepotentiallaw}
The assigned autoformalize agent should translate this physics problem into Lean declarations in the covered file, with theorem and lemma proof bodies written as `by sorry`.
\end{theorem}
\begin{proof}
This is an autoformalization task, not a proof task. Produce statements that can later be proved by the physics prover without weakening the physical model.
\end{proof}

% --- Archon declaration topology begin ---
\section{Declaration topology}
\begin{definition}[electric Potential Dimension]
  \label{def:physics:phyx-mini-0857:phyxminiproblems-problemphyxmini0857-electricpotentialdimension}
  \lean{PhyXMiniProblems.ProblemPhyXMini0857.electricPotentialDimension}
  The dimension M L² T⁻² C⁻¹ of electric potential
\end{definition}
\begin{proof}
  This is the declared model definition of electric Potential Dimension.
\end{proof}
\begin{definition}[Radial Distance Quantity]
  \label{def:physics:phyx-mini-0857:phyxminiproblems-problemphyxmini0857-radialdistancequantity}
  \lean{PhyXMiniProblems.ProblemPhyXMini0857.RadialDistanceQuantity}
  A nonnegative, unit-independent radial distance
\end{definition}
\begin{proof}
  This is the declared model definition of Radial Distance Quantity.
\end{proof}
\begin{definition}[Signed Charge Quantity]
  \label{def:physics:phyx-mini-0857:phyxminiproblems-problemphyxmini0857-signedchargequantity}
  \lean{PhyXMiniProblems.ProblemPhyXMini0857.SignedChargeQuantity}
  A signed, unit-independent electric charge
\end{definition}
\begin{proof}
  This is the declared model definition of Signed Charge Quantity.
\end{proof}
\begin{definition}[Electric Potential Quantity]
  \label{def:physics:phyx-mini-0857:phyxminiproblems-problemphyxmini0857-electricpotentialquantity}
  \lean{PhyXMiniProblems.ProblemPhyXMini0857.ElectricPotentialQuantity}
  \uses{def:physics:phyx-mini-0857:phyxminiproblems-problemphyxmini0857-electricpotentialdimension}
  A signed, unit-independent electric potential
\end{definition}
\begin{proof}
  This is the declared model definition of Electric Potential Quantity.
\end{proof}
\begin{definition}[radial Distance In Meters]
  \label{def:physics:phyx-mini-0857:phyxminiproblems-problemphyxmini0857-radialdistanceinmeters}
  \lean{PhyXMiniProblems.ProblemPhyXMini0857.radialDistanceInMeters}
  \uses{def:physics:phyx-mini-0857:phyxminiproblems-problemphyxmini0857-radialdistancequantity}
  Coherent-SI readout of a radial distance in metres
\end{definition}
\begin{proof}
  This is the declared model definition of radial Distance In Meters.
\end{proof}
\begin{definition}[radial Distance In Centimeters]
  \label{def:physics:phyx-mini-0857:phyxminiproblems-problemphyxmini0857-radialdistanceincentimeters}
  \lean{PhyXMiniProblems.ProblemPhyXMini0857.radialDistanceInCentimeters}
  \uses{def:physics:phyx-mini-0857:phyxminiproblems-problemphyxmini0857-radialdistancequantity, def:physics:phyx-mini-0857:phyxminiproblems-problemphyxmini0857-radialdistanceinmeters}
  Centimetre readout used by the radius labels in the image
\end{definition}
\begin{proof}
  This is the declared model definition of radial Distance In Centimeters.
\end{proof}
\begin{definition}[charge In Coulombs]
  \label{def:physics:phyx-mini-0857:phyxminiproblems-problemphyxmini0857-chargeincoulombs}
  \lean{PhyXMiniProblems.ProblemPhyXMini0857.chargeInCoulombs}
  \uses{def:physics:phyx-mini-0857:phyxminiproblems-problemphyxmini0857-signedchargequantity}
  Coherent-SI readout of a signed electric charge in coulombs
\end{definition}
\begin{proof}
  This is the declared model definition of charge In Coulombs.
\end{proof}
\begin{definition}[charge In Nanocoulombs]
  \label{def:physics:phyx-mini-0857:phyxminiproblems-problemphyxmini0857-chargeinnanocoulombs}
  \lean{PhyXMiniProblems.ProblemPhyXMini0857.chargeInNanocoulombs}
  \uses{def:physics:phyx-mini-0857:phyxminiproblems-problemphyxmini0857-signedchargequantity, def:physics:phyx-mini-0857:phyxminiproblems-problemphyxmini0857-chargeincoulombs}
  Nanocoulomb readout used by the source label in the image
\end{definition}
\begin{proof}
  This is the declared model definition of charge In Nanocoulombs.
\end{proof}
\begin{definition}[electric Potential In Volts]
  \label{def:physics:phyx-mini-0857:phyxminiproblems-problemphyxmini0857-electricpotentialinvolts}
  \lean{PhyXMiniProblems.ProblemPhyXMini0857.electricPotentialInVolts}
  \uses{def:physics:phyx-mini-0857:phyxminiproblems-problemphyxmini0857-electricpotentialquantity}
  Coherent-SI readout of electric potential in volts
\end{definition}
\begin{proof}
  This is the declared model definition of electric Potential In Volts.
\end{proof}
\begin{definition}[Diagram Point]
  \label{def:physics:phyx-mini-0857:phyxminiproblems-problemphyxmini0857-diagrampoint}
  \lean{PhyXMiniProblems.ProblemPhyXMini0857.DiagramPoint}
  The three observation points named in the supplied figure
\end{definition}
\begin{proof}
  This is the declared model definition of Diagram Point.
\end{proof}
\begin{definition}[Radial Circle]
  \label{def:physics:phyx-mini-0857:phyxminiproblems-problemphyxmini0857-radialcircle}
  \lean{PhyXMiniProblems.ProblemPhyXMini0857.RadialCircle}
  The two dashed circles centered on the source charge
\end{definition}
\begin{proof}
  This is the declared model definition of Radial Circle.
\end{proof}
\begin{definition}[Figure Charge Sign]
  \label{def:physics:phyx-mini-0857:phyxminiproblems-problemphyxmini0857-figurechargesign}
  \lean{PhyXMiniProblems.ProblemPhyXMini0857.FigureChargeSign}
  The sign glyph drawn inside the central source circle
\end{definition}
\begin{proof}
  This is the declared model definition of Figure Charge Sign.
\end{proof}
\begin{definition}[expected Circle For Point]
  \label{def:physics:phyx-mini-0857:phyxminiproblems-problemphyxmini0857-expectedcircleforpoint}
  \lean{PhyXMiniProblems.ProblemPhyXMini0857.expectedCircleForPoint}
  \uses{def:physics:phyx-mini-0857:phyxminiproblems-problemphyxmini0857-diagrampoint, def:physics:phyx-mini-0857:phyxminiproblems-problemphyxmini0857-radialcircle}
  Circle on which each named point appears in the primary image
\end{definition}
\begin{proof}
  This is the declared model definition of expected Circle For Point.
\end{proof}
\begin{definition}[Point Charge Potential Figure]
  \label{def:physics:phyx-mini-0857:phyxminiproblems-problemphyxmini0857-pointchargepotentialfigure}
  \lean{PhyXMiniProblems.ProblemPhyXMini0857.PointChargePotentialFigure}
  \uses{def:physics:phyx-mini-0857:phyxminiproblems-problemphyxmini0857-diagrampoint, def:physics:phyx-mini-0857:phyxminiproblems-problemphyxmini0857-radialcircle, def:physics:phyx-mini-0857:phyxminiproblems-problemphyxmini0857-figurechargesign}
  Literal presentation data transcribed from image 857.png
\end{definition}
\begin{proof}
  This is the declared model definition of Point Charge Potential Figure.
\end{proof}
\begin{definition}[Potential Reference]
  \label{def:physics:phyx-mini-0857:phyxminiproblems-problemphyxmini0857-potentialreference}
  \lean{PhyXMiniProblems.ProblemPhyXMini0857.PotentialReference}
  The convention used to choose the additive constant of the potential
\end{definition}
\begin{proof}
  This is the declared model definition of Potential Reference.
\end{proof}
\begin{definition}[Point Charge Potential Setup]
  \label{def:physics:phyx-mini-0857:phyxminiproblems-problemphyxmini0857-pointchargepotentialsetup}
  \lean{PhyXMiniProblems.ProblemPhyXMini0857.PointChargePotentialSetup}
  \uses{def:physics:phyx-mini-0857:phyxminiproblems-problemphyxmini0857-radialdistancequantity, def:physics:phyx-mini-0857:phyxminiproblems-problemphyxmini0857-signedchargequantity, def:physics:phyx-mini-0857:phyxminiproblems-problemphyxmini0857-electricpotentialquantity, def:physics:phyx-mini-0857:phyxminiproblems-problemphyxmini0857-diagrampoint, def:physics:phyx-mini-0857:phyxminiproblems-problemphyxmini0857-pointchargepotentialfigure, def:physics:phyx-mini-0857:phyxminiproblems-problemphyxmini0857-potentialreference}
  Independent physical quantities in the point-charge setup. In particular, potentialAt is an observable field and is not defined from an answer choice or from the desired 900 V value
\end{definition}
\begin{proof}
  This is the declared model definition of Point Charge Potential Setup.
\end{proof}
\begin{definition}[Matches Supplied Point Charge Figure]
  \label{def:physics:phyx-mini-0857:phyxminiproblems-problemphyxmini0857-matchessuppliedpointchargefigure}
  \lean{PhyXMiniProblems.ProblemPhyXMini0857.MatchesSuppliedPointChargeFigure}
  \uses{def:physics:phyx-mini-0857:phyxminiproblems-problemphyxmini0857-radialdistanceincentimeters, def:physics:phyx-mini-0857:phyxminiproblems-problemphyxmini0857-chargeinnanocoulombs, def:physics:phyx-mini-0857:phyxminiproblems-problemphyxmini0857-expectedcircleforpoint, def:physics:phyx-mini-0857:phyxminiproblems-problemphyxmini0857-pointchargepotentialsetup}
  The labels and geometry read directly from the primary image, together with their association to the independent physical charge and distances. No electric-potential value occurs in this data structure
\end{definition}
\begin{proof}
  This is the declared model definition of Matches Supplied Point Charge Figure.
\end{proof}
\begin{definition}[Uses School Coulomb Constant]
  \label{def:physics:phyx-mini-0857:phyxminiproblems-problemphyxmini0857-usesschoolcoulombconstant}
  \lean{PhyXMiniProblems.ProblemPhyXMini0857.UsesSchoolCoulombConstant}
  \uses{def:physics:phyx-mini-0857:phyxminiproblems-problemphyxmini0857-pointchargepotentialsetup}
  The rounded Coulomb constant used by the multiple-choice calculation. The Physlib declaration Electromagnetism.EMSystem.coulombConstant is the scalar coherent-SI constant 1 / (4 π ε₀) appearing in the governing law
\end{definition}
\begin{proof}
  This is the declared model definition of Uses School Coulomb Constant.
\end{proof}
\begin{definition}[Has Physical Point Charge Parameters]
  \label{def:physics:phyx-mini-0857:phyxminiproblems-problemphyxmini0857-hasphysicalpointchargeparameters}
  \lean{PhyXMiniProblems.ProblemPhyXMini0857.HasPhysicalPointChargeParameters}
  \uses{def:physics:phyx-mini-0857:phyxminiproblems-problemphyxmini0857-radialdistanceinmeters, def:physics:phyx-mini-0857:phyxminiproblems-problemphyxmini0857-chargeincoulombs, def:physics:phyx-mini-0857:phyxminiproblems-problemphyxmini0857-pointchargepotentialsetup}
  Positivity and separation conditions selecting the physical branch
\end{definition}
\begin{proof}
  This is the declared model definition of Has Physical Point Charge Parameters.
\end{proof}
\begin{definition}[Satisfies Point Charge Potential Law]
  \label{def:physics:phyx-mini-0857:phyxminiproblems-problemphyxmini0857-satisfiespointchargepotentiallaw}
  \lean{PhyXMiniProblems.ProblemPhyXMini0857.SatisfiesPointChargePotentialLaw}
  \uses{def:physics:phyx-mini-0857:phyxminiproblems-problemphyxmini0857-radialdistanceinmeters, def:physics:phyx-mini-0857:phyxminiproblems-problemphyxmini0857-chargeincoulombs, def:physics:phyx-mini-0857:phyxminiproblems-problemphyxmini0857-electricpotentialinvolts, def:physics:phyx-mini-0857:phyxminiproblems-problemphyxmini0857-pointchargepotentialsetup}
  For a point charge q, with the potential zero fixed at infinity, the potential at every observation point a nonzero distance r from the source is k q / r. This is a general governing law and contains no numerical target
\end{definition}
\begin{proof}
  This is the declared model definition of Satisfies Point Charge Potential Law.
\end{proof}
\begin{definition}[Answer Choice]
  \label{def:physics:phyx-mini-0857:phyxminiproblems-problemphyxmini0857-answerchoice}
  \lean{PhyXMiniProblems.ProblemPhyXMini0857.AnswerChoice}
  Labels of the four displayed answer choices
\end{definition}
\begin{proof}
  This is the declared model definition of Answer Choice.
\end{proof}
\begin{definition}[voltage In Volts]
  \label{def:physics:phyx-mini-0857:phyxminiproblems-problemphyxmini0857-answerchoice-voltageinvolts}
  \lean{PhyXMiniProblems.ProblemPhyXMini0857.AnswerChoice.voltageInVolts}
  \uses{def:physics:phyx-mini-0857:phyxminiproblems-problemphyxmini0857-answerchoice}
  Voltage in volts printed beside each displayed answer label
\end{definition}
\begin{proof}
  This is the declared model definition of voltage In Volts.
\end{proof}
\begin{definition}[recorded Dataset Answer]
  \label{def:physics:phyx-mini-0857:phyxminiproblems-problemphyxmini0857-recordeddatasetanswer}
  \lean{PhyXMiniProblems.ProblemPhyXMini0857.recordedDatasetAnswer}
  \uses{def:physics:phyx-mini-0857:phyxminiproblems-problemphyxmini0857-answerchoice}
  The answer label recorded by the supplied dataset metadata
\end{definition}
\begin{proof}
  This is the declared model definition of recorded Dataset Answer.
\end{proof}
\begin{lemma}[point C distance readout]
  \label{lem:physics:phyx-mini-0857:phyxminiproblems-problemphyxmini0857-pointc-distance-readout}
  \lean{PhyXMiniProblems.ProblemPhyXMini0857.pointC_distance_readout}
  \uses{def:physics:phyx-mini-0857:phyxminiproblems-problemphyxmini0857-radialdistanceincentimeters, def:physics:phyx-mini-0857:phyxminiproblems-problemphyxmini0857-pointchargepotentialsetup, def:physics:phyx-mini-0857:phyxminiproblems-problemphyxmini0857-matchessuppliedpointchargefigure}
  The image data place point C at a radial distance of 2.0 cm
\end{lemma}
\begin{proof}
  The conclusion follows from the governing relations and figure readouts named in the statement of point C distance readout.
\end{proof}
\begin{lemma}[potential At Point C eq coulomb Expression]
  \label{lem:physics:phyx-mini-0857:phyxminiproblems-problemphyxmini0857-potentialatpointc-eq-coulombexpression}
  \lean{PhyXMiniProblems.ProblemPhyXMini0857.potentialAtPointC_eq_coulombExpression}
  \uses{def:physics:phyx-mini-0857:phyxminiproblems-problemphyxmini0857-radialdistanceinmeters, def:physics:phyx-mini-0857:phyxminiproblems-problemphyxmini0857-chargeincoulombs, def:physics:phyx-mini-0857:phyxminiproblems-problemphyxmini0857-electricpotentialinvolts, def:physics:phyx-mini-0857:phyxminiproblems-problemphyxmini0857-pointchargepotentialsetup, def:physics:phyx-mini-0857:phyxminiproblems-problemphyxmini0857-hasphysicalpointchargeparameters, def:physics:phyx-mini-0857:phyxminiproblems-problemphyxmini0857-satisfiespointchargepotentiallaw}
  Specializing the governing point-charge law at point C gives the Coulomb expression whose independent numerical inputs are supplied above
\end{lemma}
\begin{proof}
  The conclusion follows from the governing relations and figure readouts named in the statement of potential At Point C eq coulomb Expression.
\end{proof}
% --- Archon declaration topology end ---
% --- Archon physics formalization source end ---
